% =====================================================================
%  A Training-Free Quantum-Walk Scanner for Cryptic Allosteric Sites
%  Global Quantum + AI Challenge 2026 - Cleveland Clinic
%  Phase I Concept Proposal
%
%  Compile with:  pdflatex main.tex   (twice, for references)
% =====================================================================
\documentclass[10pt,a4paper]{article}

\usepackage[T1]{fontenc}
\usepackage[utf8]{inputenc}
\usepackage{mathptmx}   % Times-like serif; swap for lmodern/newtxtext if available
\usepackage{microtype}
\usepackage{changepage}
\usepackage{multicol}
\usepackage{quantikz}
\usepackage[margin=2.2cm,top=2.1cm,bottom=2.0cm]{geometry}
\usepackage{amsmath,amssymb}
\usepackage{booktabs}
\usepackage{tabularx}
\usepackage{enumitem}
\usepackage{titlesec}
\usepackage[hidelinks]{hyperref}
\usepackage{caption}

% --- spacing / headings -----------------------------------------------
\setlength{\parskip}{0.3em}
\setlength{\parindent}{0pt}
\titlespacing*{\section}{0pt}{0.8em}{0.3em}
\titlespacing*{\subsection}{0pt}{0.55em}{0.2em}
\titleformat{\section}{\normalfont\large\bfseries}{\thesection.}{0.6em}{}
\titleformat{\subsection}{\normalfont\bfseries\itshape}{\thesubsection}{0.6em}{}
\captionsetup{font=footnotesize,labelfont=bf,skip=3pt}
\setlist[enumerate]{itemsep=2pt,topsep=3pt,parsep=0pt,leftmargin=1.6em}
\setlist[itemize]{itemsep=2pt,topsep=3pt,parsep=0pt,leftmargin=1.4em}

% --- convenience macros ------------------------------------------------
\newcommand{\Mhat}{\hat{M}}
\newcommand{\lead}[1]{\textbf{\textit{#1}}}

\begin{document}

% =====================================================================
\begin{center}
  {\LARGE\bfseries A Training-Free Quantum-Walk Scanner\\[0.25em]
   for Cryptic Allosteric Sites in Undruggable Proteins\par}
  \vspace{0.9em}
  {\small Global Quantum + AI Challenge 2026 $\cdot$ Cleveland Clinic Enterprise Challenge
   $\cdot$ Phase~I Concept Proposal\par}
\end{center}
\vspace{0.6em}

% ---------------------------------------------------------------- abstract
\begin{center}\textbf{Abstract}\end{center}
\begin{adjustwidth}{1.2cm}{1.2cm}
\small
Over 85\% of disease-causing proteins are considered undruggable, and for these the only viable
strategy is allostery. We propose a training-free scanner that takes an apo structure and
blind-predicts distal regulatory pockets. Contact topology and evolutionary conservation are encoded
as a sparse Hermitian Hamiltonian, and signal propagation is the circuit $e^{-iHt}$ acting on an
active-site source state. The readout is the average mixing matrix of that evolution: a deterministic
spectral object, doubly stochastic, with no free time parameter and no training, which is itself the
$N\times N$ connectivity matrix the challenge requires. Residues are ranked by degree-deconfounded
connectivity to the active site, standardised against background, which is the statistic the
challenge uses for scoring. An active-site-anchored Krylov projection reduces the walk to a
one-dimensional chain that preserves source-to-residue transport exactly, decoupling qubit count from
protein length and bounding two-qubit depth independently of system size. Because no parameters are
fitted, all four targets serve purely as blind tests, and a Lindblad dephasing sweep places the
classical diffusive baseline at an endpoint of the same model.
\end{adjustwidth}
\vspace{0.5em}

% =====================================================================
\section{Understanding of the Problem}

The stated need is a scanner placed upstream of screening: it ingests static structures of targets
already shelved as undruggable and returns a probability map of cryptic and allosteric sites, so
library design can be aimed at regions mechanistically connected to disease function.

Three properties of the problem constrain any solution. The pockets are cryptic, so scanning for
surface concavity misses them by construction, which is why we do not use pocket geometry as a
filter. The signal is dynamic: a site is allosteric because perturbation there reaches the active
site, so the method must model propagation. And molecular-dynamics trajectories are excluded, ruling
out the methods that currently dominate cryptic-pocket discovery.

The elastic-network hypothesis makes the problem tractable. If contact topology drives collective
dynamics, a protein reduces to a graph and atomic force fields can be set aside; the Gaussian Network
Model reproduces crystallographic B-factors from a single-parameter C$\alpha$ contact
network~\cite{bahar1997}, and topology-driven signal propagation is well
established~\cite{chennubhotla2007}. A graph is also the natural setting for a quantum walk.

% =====================================================================
\section{Prior Work and the Gap}

\lead{Classical graph methods.} Elastic network models~\cite{bahar1997,atilgan2001} predict
fluctuations and collective modes from topology; Estrada communicability
$e^{A}$~\cite{estrada2008} measures diffusive connectivity. The strongest published result in our
exact setting is bond-to-bond propensity analysis, which recovers allosteric sites knowing only the
orthosteric site~\cite{amor2016}. These are fast, MD-free and the honest baselines; their common
feature is that propagation is diffusive.

\lead{Pocket and learned predictors.} fpocket scores geometric concavity but cannot see a pocket
closed in the apo state. The state of the art for cryptic sites is PocketMiner, a graph neural
network reaching ROC-AUC $0.87$ on 39 confirmed cryptic pockets~\cite{meller2023}. Its constraint is
supervision: labelled apo--holo data is scarce (ASBench, a few hundred
entries~\cite{asbench2015}; CASBench, 91 enzymes), limiting generalisation to
novel targets. A training-free method sidesteps this.

\lead{Quantum approaches.} Continuous-time quantum walks are mature, with a spectral theory of the
average mixing matrix~\cite{godsil2013}; open-system generalisations interpolate between quantum and
classical walks~\cite{whitfield2010}; noise-assisted transport is characterised in the photosynthetic
literature~\cite{mohseni2008,rebentrost2009}; complex edge phases yield directed
walks~\cite{lu2016}. Most directly, unitary CTQW centrality was recently applied to
residue interaction networks with an IBM hardware demonstration~\cite{mohtashim2026}.

\lead{The gap.} Existing quantum work on residue networks is centrality-based on purely topological
graphs: it asks which residues are structurally important, not which are \emph{specifically coupled
to the active site}, and leaves the coherence assumption unexamined at physiological temperature. We
differ in five respects: (i) evolutionary conservation enters as on-site energy on the diagonal,
adding functional prior without densifying $H$ or breaking Hermiticity; (ii) an optional chiral phase
from elastic-network mode direction breaks time-reversal symmetry, giving a \emph{directed}
connectivity matrix no symmetric classical measure can express; (iii) coarse-graining is anchored on
the active site by a Krylov projection preserving source transport exactly; (iv) the deliverable is a
deterministic, training-free spectral object; and (v) the coherent fraction is swept, not fixed.

% =====================================================================
\section{Technical Approach}

\subsection{Hamiltonian encoding}

Nodes are residues, which keeps the output at the residue resolution the deliverable requires. $H$
is sparse and Hermitian:
\begin{equation}
H_{ij} = w_{ij}\,e^{i\varphi_{ij}} \quad (i \neq j), \qquad H_{ii} = c_i .
\end{equation}
Three physically distinct channels are fused, each entering a different part of $H$ so that none
destroys the sparsity or the Hermiticity of the others.

\emph{Off-diagonal magnitude $w_{ij}$ --- mechanical topology.} Nodes are residues represented by
C$\alpha$; an edge exists where the C$\alpha$--C$\alpha$ distance falls within $7.3$~\AA, optionally
weighted by heavy-atom contact counts at $4.5$~\AA. This is the contact network underlying the
Gaussian Network Model, whose Kirchhoff matrix is $\Gamma = D - A$. We borrow the \emph{contact
graph} only, with no GNM mode decomposition: the dynamics is
generated by $H$, not by $\Gamma$'s slow modes. The graph is naturally sparse (degree 6--10), which
bounds circuit cost downstream. Geometry beyond connectivity enters through the phase channel rather
than as a separate distance matrix, which would be largely redundant with the mechanical contacts.

\emph{Diagonal $c_i$ --- evolutionary prior.} Per-residue conservation from ESM-2
enters as on-site energy. Carrying the functional prior on the diagonal injects sequence information
without adding edges, so $H$ stays sparse and Hermitian and the circuit cost is unchanged. This
channel separates the method from purely topological prior work, and is accompanied by an ablation
against the topology-only configuration.

\emph{Phase $\varphi_{ij}$ --- directionality (optional).} The projection of the lowest
anisotropic-network-model mode along edge $(i,j)$~\cite{atilgan2001} supplies a complex phase; ANM is
an analytic eigenproblem of the elastic Hessian, not a sampling procedure, so this stays within the
no-MD constraint. Channels are z-scored before fusion and asymmetric inputs symmetrised.

Hermiticity holds channel by channel: $c_i$ real gives $H_{ii} = \overline{H_{ii}}$; contacts are
symmetrised as $(A + A^{\mathsf T})/2$; antisymmetric phases $\varphi_{ji} = -\varphi_{ij}$ give
$H_{ji} = \overline{H_{ij}}$. With phases $H$ is complex Hermitian but non-symmetric: still a valid
walk, yet capable of directed connectivity~\cite{lu2016}. Since only phase around cycles is
physical, we gauge-fix and assert $\lVert H - H^{\dagger}\rVert_F < 10^{-10}$ as a pre-flight check.
All inputs derive from apo coordinates and sequence, and the anisotropic modes are an analytic
eigenproblem rather than MD sampling, so the no-MD constraint is satisfied.

\subsection{Readout: the average mixing matrix}

The propagator is $U(t) = e^{-iHt}$. Its long-time average mixing matrix is a deterministic spectral
object~\cite{godsil2013}:
\begin{equation}
\Mhat_{ij}
  = \lim_{T\to\infty}\frac{1}{T}\int_{0}^{T}\bigl|\langle j|e^{-iHt}|i\rangle\bigr|^{2}\,dt
  = \Bigl(\sum_{r} E_r \circ E_r\Bigr)_{ij},
\end{equation}
where $E_r$ are the spectral idempotents of $H$ and $\circ$ is the Schur product. $\Mhat$ is doubly
stochastic, positive semidefinite, carries no free time parameter and requires no training. It is
delivered directly as the $N\times N$ connectivity matrix. Its quantum content is explicit: $\Mhat$
is fixed by coherent interference between eigenspaces, and is mathematically distinct from classical
communicability $e^{A}$~\cite{estrada2008} on the same graph.

Two numerical points bear on the readout. Averaging over all $t$ can suppress the
constructive-interference peaks that motivate the walk, and the idempotent grouping in~(2) is
ill-conditioned under near-degenerate spectra. We therefore benchmark $\Mhat$ against a time-resolved
score on the same graph and report a degeneracy-tolerance scan alongside it, so the choice of readout
rests on measured discrimination.

\subsection{Ranking metric}

Because $\Mhat$ is doubly stochastic, raw connectivity is confounded by degree: hub residues score
highly by construction. The ranking metric must therefore express \emph{specific} coupling to the
active site $S$ relative to background:
\begin{equation}
s(j) = \sum_{i \in S} w_i \Mhat_{ij}, \qquad
\rho(j) = \frac{s(j)}{\sum_{k} \Mhat_{kj}}, \qquad
z(j) = \frac{\rho(j) - \mu_{\mathrm{bg}}}{\sigma_{\mathrm{bg}}},
\end{equation}
with $\mu_{\mathrm{bg}}$ and $\sigma_{\mathrm{bg}}$ taken from background residues (non-active,
non-neighbour surface residues and non-functional pockets). The primary ranking statistic $z(j)$ is
by construction the same quantity the challenge's primary objective tests, which removes any
mismatch between what we optimise and what we are scored on, and avoids the degenerate solution of
ranking hubs.

\lead{Hit selection.} We exclude $S$ and its first neighbours, enforce a distal constraint
(C$\alpha$--C$\alpha$ $> 12$~\AA\ from $S$), then take the five highest $z(j)$, merging spatially
adjacent residues into one hit. Pocket-ness is not used as a filter: cryptic pockets are closed in apo structures and fpocket will
miss them, so a hard druggability gate would defeat the purpose of the task. For targets with formed pockets we cluster high-$z$ residues and annotate pocket-ness as
a druggability flag; for c-Myc, disordered and without a static pocket, the annotation is off.

\subsection{Active-site-anchored dimensionality reduction}

Ranking requires only the source-to-all transport vector $\langle j|e^{-iHt}|\psi_S\rangle$, and
that dynamics is spanned \emph{exactly} by the Krylov subspace $\mathcal{K}_K(H,\psi_S) =
\operatorname{span}\{|\psi_S\rangle, H|\psi_S\rangle, \dots, H^{K-1}|\psi_S\rangle\}$ generated from
the normalised source state $|\psi_S\rangle = \sum_{i \in S} a_i |i\rangle$. Lanczos iteration
orthogonalises it via
\begin{equation}
\beta_r |q_{r+1}\rangle = H|q_r\rangle - \alpha_r |q_r\rangle - \beta_{r-1}|q_{r-1}\rangle,
\qquad \alpha_r = \langle q_r | H | q_r \rangle,
\end{equation}
giving $Q = [q_1 \cdots q_K]$ with $q_1 = |\psi_S\rangle$ and a tridiagonal projection
$H_K = Q^{\dagger} H Q$ carrying $\alpha_r$ on the diagonal and $\beta_r$ off it. Since
$Q^{\dagger}|\psi_S\rangle = e_1$,
\begin{equation}
\langle j|e^{-iHt}|\psi_S\rangle = [Q]_{j,:}\; e^{-iH_K t}\; e_1 ,
\end{equation}
so source transport is preserved exactly in $K$ dimensions. For the quantity that determines the
ranking, retention under coarse-graining is therefore an algebraic identity, not an empirical overlap
measure, which is what the challenge asks to be demonstrated. Generic spectral clustering
preserves only generic modes and destroys within-cluster residue resolution~\cite{novo2015}. We use $K = 32$--$64$ for analysis and
$K = 8$--$16$ on hardware, decoupling qubit count from protein length. The delivered $N \times N$
$\Mhat$ is computed by exact eigendecomposition at full resolution (single-particle walks live in $N$
dimensions, so this is polynomial and trivial for $N \approx 300$--$1000$); the Krylov route supplies
the hardware track, and agreement between the two is a required consistency check.

\subsection{Back-mapping the reduction to the $N\times N$ matrix}

The reduction runs the walk in $K$ dimensions, whereas the deliverable is the $N\times N$ residue
matrix; the Lanczos isometry $Q$, whose columns already live in residue space, connects them.
Diagonalising the small tridiagonal projection $H_K = V\Theta V^{\dagger}$
($\Theta = \operatorname{diag}(\theta_r)$, an $O(K^2)$ step) and lifting each Ritz vector to
\begin{equation}
|y_r\rangle = Q v_r \in \mathbb{C}^{N}, \qquad Y = QV \ (N\times K),
\end{equation}
yields $K$ approximate eigenpairs $(\theta_r, y_r)$ of $H$ that reproduce exactly the spectral content
reachable from $|\psi_S\rangle$. The propagator restricted to the subspace is
$\tilde U(t) = Q\,e^{-iH_K t}\,Q^{\dagger} = Y\,e^{-i\Theta t}\,Y^{\dagger}$; time-averaging
$|\tilde U_{ji}(t)|^{2}$ collapses the cross terms at non-degenerate Ritz values, so the connectivity
matrix is a Gram product of squared Ritz amplitudes,
\begin{equation}
\Mhat_{ij} = \sum_{r=1}^{K} |y_r(i)|^{2}\,|y_r(j)|^{2} = (W W^{\mathsf T})_{ij}, \qquad
W_{ir} \equiv |y_r(i)|^{2},
\end{equation}
a single $N\times K$ matrix and one product, symmetric and positive semidefinite by construction.

\lead{Why this is the natural decode.} As $K \to N$ the Ritz pairs become exact eigenpairs and
$W W^{\mathsf T}$ becomes the exact average mixing matrix of Section~3.2, so the full-resolution
deliverable and the hardware-track reconstruction are one estimator at two truncations, not two
methods. Row sums equal the Krylov-projector diagonal,
$\sum_j \Mhat_{ij} = \sum_r |y_r(i)|^{2} = (Q Q^{\dagger})_{ii} \le 1$, and the deficit
$1 - \sum_j \Mhat_{ij}$ is the fraction of residue $i$'s dynamics lying outside the captured
subspace: a per-residue coverage diagnostic reported alongside the coarse-graining fidelity of
Section~5. The source column $\Mhat[\cdot, S]$ that the ranking consumes is reconstructed exactly,
error appearing only for residue pairs spectrally distant from the active site; block or restarted
Lanczos with a few probe sources tightens those off-source blocks through the same lift $Y = QV$. On
hardware only the $K$-dimensional evolution is executed; $Q$ is a classical $N\times K$ matrix, so
back-mapping is post-processing and adds nothing to the qubit budget.

\subsection{Coherence sweep}

Interference here is a modelling choice, not a claim about physiological coherence. Embedding $H$ in
a quantum stochastic walk~\cite{whitfield2010} weights coherent and incoherent channels by a single
parameter $p \in [0,1]$:
\begin{equation}
\frac{d\rho}{dt}
 = -\,i(1-p)\,[H,\rho]
 \;+\; p\sum_{i,j}\Bigl(L_{ij}\,\rho\,L_{ij}^{\dagger}
        - \tfrac{1}{2}\bigl\{L_{ij}^{\dagger}L_{ij},\,\rho\bigr\}\Bigr),
\end{equation}
with jump operators $L_{ij} = \sqrt{\gamma_{ij}}\,|i\rangle\langle j|$ acting only along existing
edges, so both channels share the same topology. At $p \to 0$ only the von Neumann term survives and
the dynamics is the coherent walk $e^{-iHt}$; at $p \to 1$ the diagonal elements obey the classical
master equation with generator $L$, giving $e^{-Lt}$. Sweeping $p$ serves three purposes at once: it is the \emph{ablation}, since the classical
control is our own endpoint on an identical graph and generator, not a separate model; it is
the \emph{noise-resilience study}, since environment-assisted transport predicts an interior optimum
on disordered graphs~\cite{rebentrost2009}; and it answers the challenge's observation that
classical approximations are too linear. We additionally report
$\Delta(j) = \Mhat[S,j] - e^{A}[S,j]$ to isolate the interference contribution against classical
communicability.

\subsection{Learning components and quantum advantage}

The predictor is training-free, its principal robustness argument. Learning appears only where it
cannot leak: ESM-2 supplies the conservation diagonal, ablated against the topology-only
configuration; Bayesian optimisation sets the two to four thresholds by leave-one-protein-out on a
disjoint calibration set; and a language-model layer turns the ranked output into a structured brief,
so chemists can use the tool without quantum training.

Single-particle walks are classically polynomial, and we claim no speed-up: predictions come from
exact classical simulation, which the challenge permits as a quantum-inspired approach. The quantum content is the physics of the model, coherent interference
evidenced by discriminative gain over $e^{A}$, plus a credible near-term hardware path. Genuine
hardware advantage arises only for multiple simultaneous excitations under exchange coupling, where
the state space grows as $\binom{N}{m}$ ($\binom{100}{3} = 161{,}700$, $\binom{100}{5} =
75{,}287{,}520$), beyond open-system classical simulation; that regime is a Phase~II stretch goal.

\subsection{Circuit specification and resource estimate}

$H_K$ is tridiagonal, hence a one-dimensional open chain with on-site energies $\alpha_r$ and
nearest-neighbour hoppings $\beta_r$. The Krylov step therefore maps an arbitrary protein contact
topology onto the structure near-term hardware executes best: no long-range couplings, no SWAP overhead, full parallelisation.

We use a unary encoding: $K$ Krylov basis states map to $K$ qubits, the working space being the
Hamming-weight-one sector. Figure~1 gives the resulting circuit. Because the chain is
one-dimensional, two odd/even layers exhaust every bond, so \emph{two-qubit depth is independent of
$K$} and set only by the Trotter number $r$, which answers the coherence-time constraint directly. Gate counts follow immediately: $2(K-1)r$ two-qubit gates at a two-qubit depth of $4r$.

\begin{figure}[t]
\centering
\scalebox{0.92}{\begin{quantikz}[column sep=4pt,row sep=5pt,transparent]
\lstick{$q_1$} & \gate{X} & \gate{R_z(\alpha_1)}\gategroup[6,steps=3,style={dashed,rounded corners,inner sep=3pt}]{} & \gate[2]{\mathrm{XY}(\beta_1)} & \qw & \meter{} \\
\lstick{$q_2$} & \qw & \gate{R_z(\alpha_2)} & & \gate[2]{\mathrm{XY}(\beta_2)} & \meter{} \\
\lstick{$q_3$} & \qw & \gate{R_z(\alpha_3)} & \gate[2]{\mathrm{XY}(\beta_3)} & & \meter{} \\
\lstick{$q_4$} & \qw & \gate{R_z(\alpha_4)} & & \gate[2]{\mathrm{XY}(\beta_4)} & \meter{} \\
\lstick{$q_5$} & \qw & \gate{R_z(\alpha_5)} & \gate[2]{\mathrm{XY}(\beta_5)} & & \meter{} \\
\lstick{$q_6$} & \qw & \gate{R_z(\alpha_6)} & & \qw & \meter{}
\end{quantikz}}
\caption{The propagation circuit at $K=6$; production runs use $K=8$--$16$ with identical structure.
The dashed block is one Trotter step, repeated $r$ times. State preparation is the single $X$ gate,
since $|\psi_S\rangle$ is $e_1$ in the Krylov basis; each step applies the diagonal $\alpha_r$ as
parallel $R_z$ rotations, then the off-diagonal $\beta_r$ as two layers of
$\mathrm{XY}(\beta)=\exp[-i\beta(XX+YY)/2]$ on odd and even bonds, each compiling to two CNOTs plus
single-qubit rotations. Measurement is in the computational basis; shots outside Hamming weight one
are discarded.}
\end{figure}

\begin{table}[t]
\caption{Circuit cost and coherence budget. Duration assumes ${\sim}120$\,ns per two-qubit layer;
bare fidelity is $(0.994)^{n_{2Q}}$ at a representative superconducting two-qubit error rate. Device
values are read from Braket device properties at run time.}
\footnotesize
\begin{tabularx}{\textwidth}{@{}lcccccX@{}}
\toprule
\textbf{Config.} & \textbf{Qubits} & \textbf{2Q gates} & \textbf{2Q depth} & \textbf{Duration} & \textbf{Bare fidelity} & \textbf{Role}\\
\midrule
$K=8$, $r=10$   & 8  & 140  & 40 & ${\sim}5\,\mu$s  & $0.43$   & Headline hardware run\\
$K=16$, $r=10$  & 16 & 300  & 40 & ${\sim}5\,\mu$s  & $0.16$   & Hardware, with post-selection\\
$K=16$, $r=20$  & 16 & 600  & 80 & ${\sim}10\,\mu$s & $0.03$   & Trotter convergence, simulator\\
$K=32$, $r=20$  & 32 & 1240 & 80 & ${\sim}10\,\mu$s & $<0.01$  & Simulator only\\
\bottomrule
\end{tabularx}
\end{table}

\lead{Reading the coherence budget.} Two limits must be separated. Circuit \emph{duration} sits
comfortably inside coherence: ${\sim}5\,\mu$s against superconducting $T_2$ of tens of microseconds
is a ratio of ${\sim}0.1$, so the walk does not decohere mid-evolution. The binding constraint is
\emph{accumulated gate error}: 300 two-qubit gates at $99.4\%$ leave a bare success probability of
$0.16$. Hence the headline run is $K=8$ and not the largest circuit we can write, $r$ is the smallest
value passing a simulator Trotter-convergence test, and $K=32$ is declared simulator-only. Post-selection on
excitation number turns much of the residual error into \emph{discarded} shots instead of biased estimates,
which is why it is the primary mitigation and ZNE only the secondary one.

\lead{Readout, platform and budget.} Sampling 64 quasi-random evolution times gives $\Mhat[S,\cdot]$
directly, since the excited qubit labels the Krylov basis vector; results project back to residues by
$Q$. Shot counts assume 2000 shots per time point with a $1.5\times$ post-selection allowance, giving
${\sim}1\%$ absolute error on probabilities of order $1/K$. Because the chain is one-dimensional, all-to-all connectivity buys us
nothing: superconducting devices on Braket (IQM, Rigetti) are both the technical and economic fit,
with per-shot rates one to two orders of magnitude below trapped-ion alternatives, and program sets
allow the 64 time points to go as a single task. Classiq provides depth-constrained synthesis across
Trotter orders under a fixed coherence budget; we report depth before and after optimisation rather
than merely noting the tool was used. A full campaign (SV1 validation, DM1 noise modelling at
$K \le 12$, a headline run with zero-noise extrapolation, a $K$-scaling study, a second target and a
cross-platform check) totals roughly $1.5\times10^{6}$ shots, on the order of a minute of pure QPU
time; it is dominated by queueing and classical pre/post-processing, not device time.

% =====================================================================
\section{Implementation Plan}

Phase~I delivers this specification; Phase~II builds and validates it. The circuit of Section~3.8 is
the defining construction; because its single-excitation sector is exactly solvable, the same $\Mhat$
is computed classically at full resolution, serving as both deliverable and hardware reference, so no
milestone depends on device availability.

\begin{enumerate}
\item \lead{Structure and features.} Fetch apo structures, select the catalytic domain, strip waters
      and heteroatoms, build the heavy-atom contact graph, and compute per-residue conservation.
\item \lead{Hamiltonian assembly.} Fuse magnitudes, diagonal and optional phases with z-score
      normalisation; gauge-fix; assert Hermiticity.
\item \lead{Readout.} Exact eigendecomposition to $\Mhat$, with the degeneracy-tolerance scan and the
      time-resolved comparison.
\item \lead{Ranking.} Compute $s$, $\rho$ and $z$; apply distal and neighbour filters; return the top
      five residues with optional pocket annotation.
\item \lead{Coherence sweep.} Lindblad dephasing scan from $p = 0$ to $p = 1$, reported against
      $e^{A}$ and $e^{-Lt}$.
\item \lead{Validation.} Blind runs on the four targets, then extended generalisation runs, against
      all baselines in Section~5.
\item \lead{Hardware.} Execute the circuits of Section~3.8 on Braket, verifying the quantum $\Mhat$
      against the classical one under post-selection and zero-noise extrapolation.
\end{enumerate}

\begin{table}[t]
\caption{Data and tools. No training set is required.}
\small
\begin{tabularx}{\textwidth}{@{}lX@{}}
\toprule
\textbf{Purpose} & \textbf{Source or tool}\\
\midrule
Required targets & RCSB PDB: 4OBE$\to$6OIM, 1OPL$\to$5MO4, 5TBY$\to$6C1H, 1NKP (c-Myc)\\
Calibration and generalisation & ASD, ASBench, CASBench; leave-one-protein-out, disjoint from targets\\
Structure and features & BioPython, ProDy (contact graph, ANM modes), fair-esm (ESM-2 conservation)\\
Simulation and readout & NumPy / SciPy eigendecomposition for $\Mhat$; QuTiP for the Lindblad sweep\\
Circuits and hardware & Qiskit, AWS Braket SDK, Classiq (synthesis, depth optimisation), Mitiq (ZNE)\\
Baselines and output & fpocket, NetworkX, ProDy GNM centrality; PyMOL and browser viewer; LLM brief\\
\bottomrule
\end{tabularx}
\end{table}

\lead{Schedule (17 Nov 2026 to 28 Feb 2027).} Weeks 1--3, structure pipeline, Hamiltonian assembly,
classical baselines. Weeks 4--6, $\Mhat$ readout, ranking metric, degeneracy and time-resolved
comparisons. Weeks 7--8, blind runs with full statistics. Weeks 9--10, coherence sweep and
calibration. Weeks 11--12, generalisation on ASD. Week 13, hardware demonstration. Weeks 14--15,
interpretation layer, visualisation, final report.

\lead{Deliverables.} (i) The $N\times N$ connectivity matrix $\Mhat$ at full residue resolution.
(ii) Ranked top-five allosteric residues per target by $z(j)$. (iii) This methodological report,
identifying $\Mhat$ as the quantum metric and coherent interference between eigenspaces as its link
to biological signal transmission. (iv) An interactive 3D connectivity map showing the protein
coloured by $z(j)$ with predicted site and propagation paths, in PyMOL and a browser viewer. The map,
not the matrix, is the artefact a chemist acts on, and the language-model brief accompanies it.

% =====================================================================
\section{Validation Plan}

\begin{table}[t]
\caption{Blind validation targets. The apo structure is the sole input; holo defines ground truth
only after predictions are frozen.}
\small
\begin{tabularx}{\textwidth}{@{}lcclc@{}}
\toprule
\textbf{Target} & \textbf{Apo} & \textbf{Holo} & \textbf{Site to blind-predict} & \textbf{Ref.}\\
\midrule
KRAS G12C       & 4OBE & 6OIM & Cryptic Switch-II pocket (sotorasib)            & \cite{ostrem2013}\\
BCR-ABL1        & 1OPL & 5MO4 & Distal myristoyl pocket (asciminib)             & \cite{wylie2017}\\
Cardiac myosin  & 5TBY & 6C1H & Site stabilising the super-relaxed state        & \cite{anderson2018}\\
c-Myc           & 1NKP & ---  & Myc--Max interface; pocket annotation disabled  & ---\\
\bottomrule
\end{tabularx}
\end{table}

\lead{Ground truth and controls.} Positive labels come from the holo structure and are used only after predictions are frozen: after
superposition onto apo coordinates, residues with a heavy atom within $4.5$~\AA\ of an
allosteric-ligand heavy atom are labelled allosteric, with indices mapped back through the apo--holo
alignment. Two negative sets match the controls named in the challenge. Random background draws
surface residues excluding the active site, its first neighbours and the labelled site.
Non-functional pockets are fpocket pockets on the apo structure overlapping neither the active site
nor any annotated functional site, scored by the same $z(j)$. The second is the stricter test, since
it excludes the possibility that the method merely separates pocket-like from non-pocket-like
surface.

\lead{Metrics and statistical tests.} Per target we report residue-level ROC-AUC, top-5 and top-10 hit rate, enrichment at top-$k$, and the
minimum C$\alpha$ distance from each predicted residue to the true site centre. Significance is
tested against each negative set by Mann--Whitney $U$ and by a permutation test resampling residues
matched on solvent accessibility and distance from the active site, so a geometric confound cannot
produce a significant result on its own. We report $p$-values with rank-biserial effect sizes and
bootstrap intervals on ROC-AUC. A target passes when both tests reach $p<0.05$ after Holm correction
across targets and the true site appears in the top five.

\lead{Baselines and ablations.} Baselines run on the identical graph wherever possible: classical communicability
$e^{A}$~\cite{estrada2008}, classical diffusion $e^{-Lt}$, GNM centrality, bond-to-bond
propensity~\cite{amor2016} as the strongest published method in this setting, PocketMiner at ROC-AUC
$0.87$~\cite{meller2023} as the learned reference, and a random ranking. The two that share our graph
and generator are decisive, since they separate the contribution of interference from that of graph
structure. Ablations remove one component at a time: the conservation diagonal, the chiral phase, the
degree deconfounding in $\rho(j)$, and the background standardisation, with the readout additionally
swapped between $\Mhat$ and a time-resolved score. Each ablation is reported as a $\Delta$ROC-AUC so
the contribution of every channel is quantified.

\lead{Robustness.} Noise resilience is assessed on the metric, not only on the state. We sweep the dephasing parameter
$p$, simulate gate-error and finite-coherence models on DM1 at $K\le 12$, and report the Spearman
correlation between noisy and noiseless $z(j)$ with the top-5 overlap: what matters is whether the
ranking survives, not whether the state does. Coarse-graining fidelity uses the same two measures
between full-resolution $\Mhat$ and its Krylov reduction across $K=8,12,16,32,64$, complementing the
exactness identity of~(5) with an empirical curve showing where truncation begins to matter.

\lead{Blind protocol and reproducibility.} No parameters are fitted to data, so leakage is structurally limited. The thresholds come from
leave-one-protein-out on a calibration set disjoint from the targets; ground-truth labels are
computed in a module the predictor cannot import, enforced by unit test; and sensitivity sweeps over
every threshold are published. Code, environment specification, seeds and intermediate matrices are
released so each figure regenerates from a clean checkout.

% =====================================================================
\section{Impact, Metrics of Success and Risks}

\lead{Where the value sits.} The tool acts where a target is advanced or abandoned. Screening without
knowing where on the surface to aim is the most expensive form of ignorance in early discovery, and
shelved targets carry zero option value. A static structure becomes a ranked surface map in minutes
instead of months, and being training-free the method applies immediately to targets with no
allosteric precedent: the population supervised predictors serve worst.

\lead{Metrics of success.} Discrimination against the published $0.87$~\cite{meller2023};
interference gain over $e^{A}$ and $e^{-Lt}$ on the same graph; runtime per target; top-5 hit rate
and distance to the true site; and generalisation on untouched ASD targets.

\lead{Risks.} Time-averaging may suppress the interference peaks that motivate the walk, which is why
a time-resolved score is benchmarked alongside $\Mhat$. Near-degenerate spectra make the idempotent
grouping in~(2) numerically delicate, addressed by tolerance scanning. The conservation prior is the
largest exposure to the ab-initio-from-topology requirement, hence the parallel topology-only run.
The elastic-network assumption is weakest for c-Myc, which is nonetheless delivered as a required
target with pocket annotation disabled. Finally, interference may add nothing, in which case the
sweep returns $p^{*}=1$; that is a controlled negative result on a question the challenge itself
raises, and the classical pipeline remains a working deliverable.

% =====================================================================
\footnotesize
\setlength{\columnsep}{16pt}
\begin{multicols}{2}
\begin{thebibliography}{99}
\setlength{\itemsep}{0pt}\setlength{\parskip}{0pt}\linespread{0.95}\selectfont

\bibitem{godsil2013} Godsil C. Average mixing of continuous quantum walks.
  \textit{J. Comb. Theory Ser. A} \textbf{120}, 1649--1662 (2013). arXiv:1103.2578

\bibitem{mohtashim2026} Mohtashim SI, Sajjan M, Kais S. Continuous-time quantum-walk centrality for
  protein residue interaction networks. \textit{J. Am. Chem. Soc.} \textbf{148}, 29206--29219 (2026).
  doi:10.1021/jacs.6c08053

\bibitem{amor2016} Amor BRC, Schaub MT, Yaliraki SN, Barahona M. Prediction of allosteric sites and
  mediating interactions through bond-to-bond propensities. \textit{Nat. Commun.} \textbf{7}, 12477
  (2016). doi:10.1038/ncomms12477

\bibitem{novo2015} Novo L, Chakraborty S, Mohseni M, Neven H, Omar Y. Systematic dimensionality
  reduction for quantum walks. \textit{Sci. Rep.} \textbf{5}, 13304 (2015). doi:10.1038/srep13304

\bibitem{whitfield2010} Whitfield JD, Rodr\'iguez-Rosario CA, Aspuru-Guzik A. Quantum stochastic
  walks. \textit{Phys. Rev. A} \textbf{81}, 022323 (2010). doi:10.1103/PhysRevA.81.022323

\bibitem{mohseni2008} Mohseni M, Rebentrost P, Lloyd S, Aspuru-Guzik A. Environment-assisted quantum
  walks in photosynthetic energy transfer. \textit{J. Chem. Phys.} \textbf{129}, 174106 (2008).
  doi:10.1063/1.3002335

\bibitem{rebentrost2009} Rebentrost P, Mohseni M, Kassal I, Lloyd S, Aspuru-Guzik A.
  Environment-assisted quantum transport. \textit{New J. Phys.} \textbf{11}, 033003 (2009).
  doi:10.1088/1367-2630/11/3/033003

\bibitem{estrada2008} Estrada E, Hatano N. Communicability in complex networks.
  \textit{Phys. Rev. E} \textbf{77}, 036111 (2008). doi:10.1103/PhysRevE.77.036111

\bibitem{lu2016} Lu D, Biamonte JD, Li J, et al. Chiral quantum walks. \textit{Phys. Rev. A}
  \textbf{93}, 042302 (2016). doi:10.1103/PhysRevA.93.042302

\bibitem{bahar1997} Bahar I, Atilgan AR, Erman B. Direct evaluation of thermal fluctuations in
  proteins using a single-parameter harmonic potential. \textit{Folding \& Design} \textbf{2},
  173--181 (1997). doi:10.1016/S1359-0278(97)00024-2

\bibitem{atilgan2001} Atilgan AR, Durell SR, Jernigan RL, et al. Anisotropy of fluctuation dynamics
  of proteins with an elastic network model. \textit{Biophys. J.} \textbf{80}, 505--515 (2001).
  doi:10.1016/S0006-3495(01)76033-X

\bibitem{chennubhotla2007} Chennubhotla C, Bahar I. Signal propagation in proteins and relation to
  equilibrium fluctuations. \textit{PLoS Comput. Biol.} \textbf{3}, 1716--1726 (2007).
  doi:10.1371/journal.pcbi.0030172

\bibitem{meller2023} Meller A, Ward M, Borowsky J, et al. Predicting locations of cryptic pockets
  from single protein structures using the PocketMiner graph neural network.
  \textit{Nat. Commun.} \textbf{14}, 1177 (2023). doi:10.1038/s41467-023-36699-3

\bibitem{asbench2015} Huang Z, et al. ASBench: benchmarking sets for allosteric discovery.
  \textit{Bioinformatics} \textbf{31}, 2357--2359 (2015). doi:10.1093/bioinformatics/btv169

\bibitem{ostrem2013} Ostrem JM, Peters U, Sos ML, Wells JA, Shokat KM. K-Ras(G12C) inhibitors
  allosterically control GTP affinity and effector interactions. \textit{Nature} \textbf{503},
  548--551 (2013). doi:10.1038/nature12796

\bibitem{wylie2017} Wylie AA, Schoepfer J, Jahnke W, et al. The allosteric inhibitor ABL001 enables
  dual targeting of BCR--ABL1. \textit{Nature} \textbf{543}, 733--737 (2017). doi:10.1038/nature21702

\bibitem{anderson2018} Anderson RL, Trivedi DV, Sarkar SS, et al. Deciphering the super relaxed state
  of human $\beta$-cardiac myosin and the mode of action of mavacamten. \textit{PNAS} \textbf{115},
  E8143--E8152 (2018). doi:10.1073/pnas.1809540115

\end{thebibliography}
\end{multicols}

\end{document}
